
\subsection{Cosmological Impedance Route for the Background Density}
\label{sec:rt_rhof_cosmological_impedance_route}

\paragraph{Status.}
\textbf{[RESEARCH TRACK / NOT CANON-DERIVED].}
This subsection records a candidate route for reducing the calibrated
background density \(\rho_{\!f}\) to a cosmological vacuum input plus a
UV--IR impedance lift. It does not promote \(\rho_{\!f}\) out of the
primitive calibration set in the main canon.

The orthodox vacuum mass-equivalent density associated with a cosmological
constant is
\begin{align}
    \rho_\Lambda
    =
    \frac{\Lambda c^2}{8\pi G} .
\end{align}
For a spatially flat \(\Lambda\)CDM background,
\begin{align}
    \Lambda
    =
    \frac{3\Omega_\Lambda H_0^2}{c^2},
    \qquad
    \rho_\Lambda
    =
    \Omega_\Lambda\frac{3H_0^2}{8\pi G} .
\end{align}
Using the Planck-2018 baseline values
\(H_0\simeq 67.4\,\mathrm{km\,s^{-1}\,Mpc^{-1}}\) and
\(\Omega_\Lambda\simeq0.685\), this gives
\begin{align}
    \Lambda &\simeq 1.09\times10^{-52}\,\mathrm{m^{-2}},\\
    \rho_\Lambda &\simeq 5.85\times10^{-27}\,\mathrm{kg\,m^{-3}} .
\end{align}
The ratio required to obtain the canonical SST background density is
\begin{align}
    \frac{\rho_{\!f}}{\rho_\Lambda}
    \simeq
    1.20\times10^{20} .
\end{align}
This is close to the one-power core--Planck ratio
\begin{align}
    \frac{r_c}{\ell_P}
    =
    8.72\times10^{19},
\end{align}
which suggests the line-impedance ansatz
\begin{align}
    \rho_{\!f}
    =
    \chi_\Lambda
    \left(\frac{r_c}{\ell_P}\right)
    \rho_\Lambda,
    \qquad
    \chi_\Lambda
    \simeq
    1.37 .
    \label{eq:rt_rhof_lambda_chi}
\end{align}
Two non-equivalent simple closures are numerically relevant:
\begin{align}
    \rho_{\!f}^{(c_s)}
    &:=
    \sqrt{2}
    \left(\frac{r_c}{\ell_P}\right)
    \rho_\Lambda
    \simeq
    7.21\times10^{-7}\,\mathrm{kg\,m^{-3}},
    \\
    \rho_{\!f}^{(\Omega)}
    &:=
    2\Omega_\Lambda
    \left(\frac{r_c}{\ell_P}\right)
    \rho_\Lambda
    \simeq
    6.98\times10^{-7}\,\mathrm{kg\,m^{-3}} .
\end{align}
The first closure uses the internal core-acoustic factor already present in
\(c_s=\mathbf{v}_{\!\boldsymbol{\circlearrowleft}}/\sqrt{2}\), but does not yet
derive why the cosmological vacuum sector should couple through that
impedance. The second closure is numerically sharper for Planck-2018
parameters, but it introduces an explicitly epoch-dependent projection
\(\Omega_\Lambda\) and is therefore not acceptable as a fundamental
constant derivation unless SST supplies a time-independent interpretation of
that factor.

\paragraph{Dimensional interpretation.}
The power of \(r_c/\ell_P\) must be one, not two or three. A volume lift
\((r_c/\ell_P)^3\) or area lift \((r_c/\ell_P)^2\) overshoots catastrophically.
The only viable route is therefore a one-dimensional UV--IR mode-count or
line-impedance theorem. In physical terms, the cosmological vacuum density
would have to be lifted into a local inertial background density by counting
Planck-normal modes along one core/horn length, not across an area or volume.

\paragraph{Canonical conclusion.}
Equation~\eqref{eq:rt_rhof_lambda_chi} is a promising falsifier route, but
not a completed derivation. The main canon should continue to label
\(\rho_{\!f}\) as a calibrated effective background density. Promotion to
\textbf{[DERIVED]} requires an SST theorem that derives either
\(\chi_\Lambda=\sqrt{2}\), or an epoch-invariant replacement for
\(2\Omega_\Lambda\), from the coarse-graining map between the cosmological
vacuum sector and the local incompressible swirl medium.
